To test the predictor end to end with measurements we control, we use the only hardware available
to us: an Apple M5 laptop (16\,GB unified memory). \texttt{bottle probe} measures its peaks
(GPU via Metal/MPS: 12.1 fp16 / 2.9 fp32 TFLOPS, 121\,GB/s; CPU: 2.1 fp32 TFLOPS, 127\,GB/s),
which define three ``devices'': CPU fp32, GPU fp32 and GPU fp16. We profile all eight VLA presets
on each with the synthetic profiler (models above 1B parameters are profiled on 4 layers per stage
on the GPU and 2 on the CPU, then extrapolated per layer), and predict every device from every
other (48 predictions per calibrated predictor), including cross-precision transfer.

\begin{table}[h]\centering\small
\begin{tabular}{lrr}
\toprule
Predictor & MAPE & median APE \\
\midrule
Roofline, ideal efficiency & 36.6\% & 34.1\% \\
Roofline, default efficiency & 31.8\% & 31.6\% \\
Global calibration & 31.5\% & 27.2\% \\
Stage calibration & 29.5\% & 25.3\% \\
Overhead-aware calibration (\textbf{ours}) & 26.2\% & 21.1\% \\
\bottomrule
\end{tabular}

\caption{E1: latency prediction across the three local Apple M5 configurations (8 VLAs).}\label{tab:local}
\end{table}

Table~\ref{tab:local} shows that with probed specs the roofline alone is already within 32\% MAPE
(19.5\% on CPU, 29.2\% GPU fp16, 46.6\% GPU fp32), because the synthetic profiler runs pure
transformer stacks with little host overhead. Calibration helps modestly, and overhead-aware
calibration is best (26.2\%). The largest difference appears in cross-precision transfer: from
GPU fp16 to GPU fp32, global calibration errs by 57.3\% and overhead-aware by 24.8\%, because
the fp16 run's fixed overheads do not slow down with fp32 arithmetic. Two caveats bound this
experiment. First, layer extrapolation itself costs accuracy: SmolVLA on GPU fp16 extrapolated
from 2 layers is 27\% off its full-depth measurement. Second, the laptop is a shared, unmanaged
machine: the same full-depth SmolVLA fp16 configuration measured 104\,ms and 167\,ms in two runs
under different memory pressure, a run-to-run spread comparable to the prediction error.
E1 therefore shows that the pipeline works on a no-cost device and that overhead-aware
calibration is not worse than the alternatives when overhead is small; the decisive test is E2.
